\documentclass[book.tex]{subfiles}

\begin{document}
\chapter{Covariant Derivatives and Connections}
\label{chap:covariant_derivative}
\noindent\rule{11cm}{0.4pt} \\
\textbf{Prerequisites:} \autoref{chap:manifold} \\
\textbf{Difficulty Level:} *** \\
\noindent\rule{11cm}{0.4pt}
\vspace{5mm}

The covariant derivative provides a method of defining derivatives to be taken along the tangent vectors of a manifold. Importantly, the covariant derivative is defined so as to be independent of the choice of local coordinate system. 

\section{Mathematical Definition}

Let $M$ be a manifold and let $p \in M$ be a point on the manifold. Let $f: M \to \mathbb{R}$ be a real valued function on the manifold. Let 
\begin{align}
v : M \to T_p M
\end{align}
be a vector field on $M$. Then we denote the covariant derivative 
\begin{align}
(\nabla_v f)_p \in T_p M,
\end{align}
to be the tangent vector at $p$ such that the following properties hold

\begin{align}
(\nabla_{gx + hy} u)_p &= (\nabla_x u)_p g + (\nabla_y u)_p h \\
(\nabla_v (u + w))_p &= (\nabla_v u)_p + (\nabla_v w)_p \\
(\nabla_v (fu))_p &= f(p)(\nabla_v u)_p  + (\nabla_v f)_p u_p
\end{align}

The goal of this definition is to make the definition of the covariant derivative independent of the choice of coordinate basis. For another definition, suppose that we have a differentiable curve
\begin{align}
\phi:[-1, 1] \to M
\end{align}
such that $\phi(0) = p$ and $\phi'(0) = v$. Then we define the covariant derivative to be
$$(\nabla_v f)_p = (f \circ \phi)'(0)$$

$$= \lim_{t \to 0} \frac{f(\phi(t)) - f(p)}{t}$$
We can then view the covariant derivative as a function mapping points in the manifold $M$ to real scalars $(nabla_v f)_p$
\begin{align}
\nabla_v f: M \to \mathbb{R}
\end{align}
%TODO: Tie this to Christoffel symbols

%\section{Exercises}
%\begin{enumerate}
%    \item TODO
%\end{enumerate}

\end{document}